
\clearpage
\section*{35. 整係数関数からなる極大領域について}
\begin{center}
E. Noether\\
\textit{Rec. Soc. Math. Moscou \textbf{36} (1929), pp. 65--72}
\end{center}

以下の考察は、相対的に整な関数に関する Hilbert の問題\footnote{D. Hilbert, \emph{Mathematische Probleme} (Gott. Nachr. 1900, pp. 253--297), 第14問題。}に密接に関わるものであり、E. Hecke の折に触れた質問に由来する。Hecke は、自分が補題として必要とした一つの特殊な場合を、直接計算によって処理することができたのである。\footnote{E. Hecke, ``Über die Konstruktion relativ-Abelscher Zahlkörper durch Modulfunktionen von zwei Variabeln'' (Math. Ann. 74 (1913), pp. 465--510), \S\,2。Hecke では、有限な整域 $\mathfrak{I}$ が、二変数の冪級数の三つの商によって生成され、それらの間には一つの代数関係が成り立つ。}

問題設定は次の通りである。与えられた領域 $\mathfrak{Z}$ の内部で、$z_1,\ldots,z_r$ の整係数関数、すなわち多項式、冪級数、またはそれらの商からなる\emph{極大領域}とは、整数分母を付け加えてももはや拡大できない領域のことである。すなわち、$g(x)$ が $\mathfrak{Z}$ に属し、$a$ が非零の有理整数であるとき、$a\cdot g(x)$ の所属から常に $g(x)$ の所属が従うものをいう。$\mathfrak{Z}$ 内の任意の整域 $\mathfrak{I}$ は、一意的に極大領域 $M(\mathfrak{I})$ を生成する。これは、ある $a\ne0$ について $a\cdot g(x)$ が $\mathfrak{I}$ に属するような、$\mathfrak{Z}$ のすべての関数 $g(x)$ から成る。

問われるのは、もとの整域 $\mathfrak{I}$ が有限であった場合、すなわち $\mathfrak{Z}$ の有限個の関数 $f_1(x),\ldots,f_t(x)$ に関する整係数多項式全体から成っていた場合、この極大な $M(\mathfrak{I})$ の分母 $a$ について何が言えるかである。このとき $\mathfrak{Z}$ としては、$z_1,\ldots,z_r$ のすべての整係数多項式または冪級数、あるいは $f(x)$ に現れる分母とその冪以外の分母を含まない商を取るものとする。

私は、冪級数またはその商の場合には、$f(x)$ 相互の依存関係が代数的なものだけであるという付加仮定のもとで、$a$ を常に、そこに現れる相異なる素数が有限個に限られるように選べることを示す。ただし、一般にはそれらの素数は任意の高い冪で現れ得る。

後者は直ちに明らかである。$a\cdot g(x)$ は $\mathfrak{I}$ に属するが、どの指数 $\varrho$ に対しても $a^\varrho g(x)$ があらかじめ制限された形で $\mathfrak{I}$ に入らないならば、すべての冪 $a^\varrho$ が現れる。例えば $\mathfrak{I}=[2x]$ のとき、$x=(2x)/2$、$x^2=(2x)^2/2^2$ である。

有界性の問題は次のようにしか定式化できない。一般には無限であるような $M(\mathfrak{I})$ の生成系を、対応する分母における素数の指数が有界に留まるように与えられるか、という問題である。関与する素数は有限個であるから、この問題は既知の議論\footnote{Hilbert 前掲、または E. Noether, \emph{Die Endlichkeit des Systems der ganzzahligen Invarianten binärer Formen} (Gött. Nachr. 1919, pp. 138--156), \S\,5 のより詳しい説明を参照。}によって $M(\mathfrak{I})$ の有限性の問題と同値であり、したがって整係数の形での、相対的に整な関数に関する Hilbert の問題そのものである。この問題は、ここで用いる方法では決定できない。\footnote{したがって例えば、基本形式系の整（射影的）不変式の場合には、それらが通常の意味での完全な不変式系によって、しかも分母に有限個の相異なる素数だけを含んで表されることが示される。これは通常の $\Omega$ 過程の方法からは従わない結果である。しかし整な有限性そのものについては何も述べていない。それはこれまで二元基本形式の場合にだけ知られている（3) で引用した注を参照）。}

分母 $a$ に現れる相異なる素数 $p$ が有限個だけであることの証明は、各そのような $p$ に対し、$\mathfrak{I}$ を生成する $f(x)$ の間に少なくとも一つの法 $p$ の関係が対応する、という事実に基づく。この関係は整数係数をもち、$x$ において恒等的に満たされるものではない。別の言い方をすれば、$\mathfrak{o}$ を不定元 $u_1,\ldots,u_t$ に関する整係数多項式領域とすると、対応 $u_i\mapsto f_i(x)$ によって $\mathfrak{I}$ は $\mathfrak{o}$ の準同型像となる。この準同型は $\mathfrak{o}$ の素イデアル $\mathfrak{a}$ によって媒介される、すなわち $\mathfrak{I}$ は剰余類環 $\mathfrak{o}/\mathfrak{a}$ と同型になる。これに対応して、$\mathfrak{I}_p$ は $\mathfrak{o}_p$ の準同型像となる。添字は法 $p$ の剰余類への移行を意味し、これは $\mathfrak{I}$ の分母に現れる有限個の素数を除いて可能である。この準同型は再び $\mathfrak{o}_p$ の素イデアル $\mathfrak{c}_p$ によって媒介され、これは $\mathfrak{a}_p$ を含む。すると、$p$ が $M(\mathfrak{I})$ の分母に現れることと、$\mathfrak{c}_p$ が $\mathfrak{a}_p$ の真の上位イデアル（イデアルの意味での真の約数）になることとは同値である。これは、$\mathfrak{I}_p$ の超越次数が $\mathfrak{I}$ のそれに比べて下がるか、あるいは $\mathfrak{a}_p$ が素イデアル性を失う場合にだけ起こり得る。\footnote{系の「超越次数」とは、周知のように、その系が代数的に依存する代数的独立な関数の個数、しかも不変な個数のことである。代数的独立な関数から成る系を既約系と呼ぶ。素イデアルの超越次数（次元）とは、その剰余類体の超越次数をいう。素イデアルは同じ超越次数をもつ真の素約数をもたない。例えば B. L. v. d. Waerden, ``Zur Nullstellentheorie der Polynomideale'', Math. Ann. 96 (1926), pp. 183--208 を参照。} これが有限個の $p$ についてしか起こらないことは、法 $p$ においても代数的依存関係が関数行列式の消滅をもたらすこと（\S\,1）、さらに $\mathfrak{a}$ が絶対素イデアルであり、したがってその素イデアル性を失うのは有限個の $p$ を法とする場合だけであることから従う。\footnote{E. Noether, ``Eliminationstheorie und allgemeine Idealtheorie'', Math. Ann. 90 (1923), pp. 229--261, \S\,7。この定理は消去理論により、絶対既約多項式がその性質を失うのは有限個の素数を法とする場合に限られるというより簡単な定理から従う（A. Ostrowski, ``Zur arithmetischen Theorie der algebraischen Größen'', Gött. Nachr. 1919, pp. 279--298, p. 296。より簡単な証明は E. Noether, ``Ein algebraisches Kriterium für absolute Irreduzibilität'', Math. Ann. 85 (1922), pp. 26--33, \S\,7）。$\mathfrak{I}$ がその超越次数よりちょうど一つ多い関数を含む整基底を持つ場合、例えば Hecke が扱った場合（注2）のように、$\mathfrak{a}$ は主イデアルになるので、消去理論を避けてこのより簡単な定理を用いるだけで足りる。素数の代わりに、どこでも（有限な）代数的数体の素イデアルを置くことができる。この定理は $\mathfrak{a}$ が零イデアルの場合にも明らかに成り立つ。}

なお、すべての考察は、有理整数の代わりに（有限な）代数的数体の整数とその素イデアルを用いる場合にも、わずかな修正で保たれる（\S\,4）。

\subsection*{\S\,1. 特性 $p$ の係数領域における関数行列式}

\noindent\textbf{1.}
$f_1(x_1,\ldots,x_r),\ldots,f_t(x_1,\ldots,x_r)$ を、有理関数、あるいはより一般に、形式的に定義された冪級数またはその商とし、その係数は特性 $p$ の体 $P$ に属するとする。導関数 $\partial f_i/\partial x_k$ も形式的に定義され、通常のすべての計算法則が保たれるものとする。$\Omega$ は $P$ の代数閉拡大体を意味する。

そのとき、特性零の場合と同様に次が成り立つ。

\noindent\textbf{定理.}
\emph{$f(x)$ の間に $\Omega$ 係数の代数的依存関係が存在するなら、関数行列 $(\partial f_i/\partial x_k)$（$i=1,\ldots,t$、$k=1,\ldots,r$）の階数は $t$ より小さい。}\footnote{元の形では、係数領域 $P$ を完全体と仮定していた。$P$ から $\Omega$ へ移ることによって、任意の体、完全でない体にまで拡張できることは B. L. v. d. Waerden に負う。特性 $p$ では、$P$ が完全であっても逆は成り立たないことは、例 $f=x^p$, $\partial f/\partial x=0$ が示している。}

\noindent\emph{証明.}
多項式領域 $\Omega[u_1,\ldots,u_t]$ において、$G(f)=0$ が $x$ に関して恒等的に成り立つようなすべての多項式 $G(u)$ の素イデアルを $\mathfrak{m}$ とする。仮定により $\mathfrak{m}$ は零イデアルではない。$\mathfrak{m}$ の中で最小次数の多項式 $G(u)\ne0$（したがって定数ではない）を取る。特に $G(u)$ は分解不能である。すると通常通り
\[
\frac{\partial G}{\partial u_1}\frac{\partial f_1}{\partial x_k}+\cdots+
\frac{\partial G}{\partial u_t}\frac{\partial f_t}{\partial x_k}=0
\qquad(k=1,\ldots,r)
\]
が得られる。したがって、関数行列の階数が $t$ より小さいか、またはすべての $\partial G/\partial u_i$ が消えるかのいずれかである。しかし $G(u)$ は最小次数の多項式として選ばれていたから、後者ならすべての $\partial G/\partial u_i$ が消えることになる。

よって $G(u)=H(u^p,\ldots,u^p)$ となる。さらに係数領域 $\Omega$ は代数閉、したがって完全であるから、$G(u)=H(u)^p=(H(u))^p$ となり、$G(u)$ の選び方に反する。証明はもちろん特性零の場合にも通用し、その場合には最後の推論が不要になる。

\medskip\noindent\textbf{2. 系.}
$F_1(x_1,\ldots,x_r),\ldots,F_t(x_1,\ldots,x_r)$ を、関数行列がちょうど階数 $t$ をもつ整係数関数（多項式、冪級数またはその商）とする。素数 $p$ は、$F(x)$ の分母にも、関数行列のすべての $t$ 行小行列式にも入らないように選ぶ。

すると $F_1(x),\ldots,F_t(x)$ は法 $p$ でも代数的独立のままである。

\noindent\emph{証明.}
$\bar P$ を法 $p$ の剰余類体とし、$f_1(x),\ldots,f_t(x)$ を、$F_1(x),\ldots,F_t(x)$ の整係数をその法 $p$ の剰余類で置き換えて得られる $\bar P$ 上の関数とする。$p$ は $F_i(x)$ の分母に入っていなかったので $f_i(x)$ は定義される。また、関数行列 $(\partial f_i/\partial x_k)$ も、$(\partial F_i/\partial x_k)$ の整係数を法 $p$ の剰余類で置き換えて得られる。$(\partial F_i/\partial x_k)$ には $F(x)$ にある以外の分母は現れない。$p$ の選び方により、$(\partial f_i/\partial x_k)$ は階数 $t$ を保つ。したがって $f(x)$ は $\Omega$ で代数的独立であり、まして $\bar P$ でも代数的独立である。

\subsection*{\S\,2. 定理の定式化}

\noindent\textbf{1.}
導入部と同様、$\mathfrak{I}$ は $z_1,\ldots,z_r$ の整係数関数、すなわち有理整数係数の多項式または冪級数、あるいはそれらの商からなる整域（零因子をもたない環）を表すものとする。冪級数の場合には、形式的に定義されたものでも、ある領域で収束するものでもよい。用いられるのは、現れる関数が整域をなすという事実だけである。

$\mathfrak{Q}$ を $\mathfrak{I}$ の商体とし、$\mathfrak{Z}$ を $\mathfrak{I}$ と $\mathfrak{Q}$ の間にある整域とする。導入部と同様に、$M$ が $\mathfrak{Z}$ において\emph{極大}であるとは、$a\ne0$ が整数で $g(x)$ が $\mathfrak{Z}$ に属するとき、$a\cdot g(x)$ の所属から常に $g(x)$ の所属が従うことをいう。$\mathfrak{I}$ に対しては、$\mathfrak{Z}$ 内で $\mathfrak{I}$ によって生成される極大領域 $M(\mathfrak{I})$ が一意的に存在する。これは $\mathfrak{Z}$ 内の $\mathfrak{I}$ を含むすべての極大領域の共通部分として得られる。なぜなら $\mathfrak{Z}$ 自身がそのような領域であり、任意個のそのような領域の共通部分も極大で、かつ $\mathfrak{I}$ を含むからである。$M(\mathfrak{I})$ は、ある $a\ne0$ について $a\cdot g(x)$ が $\mathfrak{I}$ に属するような $\mathfrak{Z}$ のすべての関数 $g(x)$ から成る。

実際、これらの $g(x)$ は、$\mathfrak{I}$ を含む $\mathfrak{Z}$ 内の極大領域 $N$ をなす。というのも、$b\cdot h(x)$ が $N$ に属し、$b\ne0$、$h(x)\in\mathfrak{Z}$ なら、ある $a\ne0$ に対して $ab\cdot h(x)$ が $\mathfrak{I}$ に属し、$ab\ne0$ であるから、$h(x)$ は $N$ に属する。しかも $\mathfrak{I}$ を含む任意の極大領域 $M$ は $N$ を含む。なぜなら $a\cdot g(x)$ が $\mathfrak{I}$ に属すれば、それは $M$ にも属し、したがって $g(x)$ も $M$ に属するからである。

\medskip\noindent\textbf{2.}
いま $\mathfrak{I}$ を、単位元をもつ有限な整域と仮定し、$f_1(x),\ldots,f_s(x)$ をその整基底とする。すなわち $\mathfrak{I}$ は $f_i(x)$ のすべての整係数多項式から成る。さらに $\mathfrak{Z}$ は、$x$ の整係数多項式または冪級数、あるいは有限個の $f_i(x)$ に現れる分母以外の分母を含まない商から成るものとする。ただし当然、それらの分母は任意の冪で現れてよい。

また、少なくとも一つの $f_i(x)$ --一般にはすべて-- に $x$ が実際に現れると仮定する。そうでなければ、$\mathfrak{I}$ と $\mathfrak{Z}$ は単に一つの分数 $1/e$ に関する整係数多項式全体から成り、証明すべきことはない。

多項式または有理関数の場合には、商体 $\mathfrak{Q}$ は有理数体 $K$ に関して超越次数 $t$（$1\le t\le r$）をもつ。例えば $f_1(x),\ldots,f_t(x)$ が既約系\footnote{「既約系」とは、周知のように、代数的独立な関数からなる系をいう。}をなすなら、$\mathfrak{Q}$ は純超越体 $K(f_1(x),\ldots,f_t(x))$ の有限代数拡大である。同時に、$f_1(x),\ldots,f_t(x)$ の関数行列は正確に階数 $t$ をもつ。\footnote{この付加仮定は、Hecke が考察した場合（注2）のように解析関数を扱い、$f_1(x),\ldots,f_t(x)$ が解析的に独立で、$f_{t+1}(x),\ldots,f_s(x)$ がこれらに代数的に依存する場合には満たされる。しかしこの仮定は意図的に形式的に、すなわち関数行列の階数として述べられている。これにより、形式的に定義された冪級数の場合にも明確な意味をもつ。したがってすべての考察は純粋に形式代数的な範囲に留まり、解析関数に関する定理は、特殊な場合にこの仮定が満たされることを確認するためだけに必要となる。}

冪級数またはその商の場合にもこの $\mathfrak{Q}$ の構造を保つには、付加仮定が必要である。すなわち、$t$ が $f_i(x)$ の関数行列の階数であり、さらに $f_1(x),\ldots,f_t(x)$ の関数行列の階数でもあるなら（したがって \S\,1.1 によりこれらは既約系をなす）、$\mathfrak{Q}$ は再び純超越体 $K(f_1(x),\ldots,f_t(x))$ の有限代数拡大であると仮定する。言い換えれば、$f_{t+1}(x),\ldots,f_s(x)$ は $K(f_1(x),\ldots,f_t(x))$ に関して代数的でなければならない。\footnote{この付加仮定は、Hecke が扱った場合（注2）のように解析関数を扱い、$f_1(x),\ldots,f_t(x)$ が解析的に独立で、$f_{t+1}(x),\ldots,f_s(x)$ がこれらに代数的に依存する場合には満たされる。} このときも $t\le r$ である。なぜなら、特性零ではすべての $\partial f_i/\partial x_k$ が同時に消えることはないからである。

同時に、この条件は、関数行列が正確に階数 $t$ をもつような整基底の任意の $t$ 個の関数系について従う。というのも、それらは既約系をなしており、残りの関数はそれらに代数的に依存するからである。

\medskip\noindent\textbf{3.}
いま $\mathfrak{o}$ を不定元 $u_1,\ldots,u_s$ に関する整係数多項式領域とする。すると対応 $u_i\mapsto f_i(x)$ が直ちに示すように、$\mathfrak{I}$ は $\mathfrak{o}$ の環準同型像である。したがって $\mathfrak{I}$ は $\mathfrak{o}/\mathfrak{a}$ と環同型である。ここで $\mathfrak{a}$ は素イデアルである。なぜなら $\mathfrak{I}$ は零因子をもたない環だからである。$\mathfrak{a}$ は、$G(f)$ が $x$ において恒等的に消えるようなすべての多項式 $G(u)$ から成る。

導入部ですでに述べたように、問題は次の定理である。

\medskip\noindent\textbf{定理.}
\emph{2 の付加仮定のもとで、$\mathfrak{o}_p$, $\mathfrak{a}_p$, $\mathfrak{I}_p$ を、$\mathfrak{o}$, $\mathfrak{a}$, $\mathfrak{I}$ から整係数を法 $p$ の剰余類で置き換えて得られる領域とする。すると、有限個を超えない例外素数を除いて --この中には、$\mathfrak{I}_p$ が定義されるために、$f(x)$ の分母に現れるものが含まれる--、$\mathfrak{o}_p$ から $\mathfrak{I}_p$ への準同型は $\mathfrak{a}_p$ によって媒介される。したがって $\mathfrak{I}_p$ は $\mathfrak{o}_p/\mathfrak{a}_p$ と、さらに $\mathfrak{o}/(\mathfrak{a},p\mathfrak{o})$ と環同型になる。}

この定理から、$M(\mathfrak{I})$ の分母 $a$ に現れる素数が有限個だけであるという事実は、次の系によって得られる。

\medskip\noindent\textbf{系.}
\emph{$\mathfrak{I}_p$ が $\mathfrak{o}_p/\mathfrak{a}_p$ と同型になるように $p$ を選ぶなら、$p\cdot g(x)$ が $\mathfrak{I}$ に属し、$g(x)$ が $\mathfrak{Z}$ に属するとき、常に $g(x)$ は $\mathfrak{I}$ に属する。}

別の言い方をすれば、$\mathfrak{I}$ は、有限個の例外素数以外のすべての素数に関して、$\mathfrak{Z}$ 内で極大である。

$\mathfrak{o}_p/\mathfrak{a}_p$ は $\mathfrak{o}/(\mathfrak{a},p\mathfrak{o})$ と同型である。なぜなら、定義により $\mathfrak{o}_p$ は剰余類環 $\mathfrak{o}/p\mathfrak{o}$ であり、対応して $\mathfrak{a}_p$ は $(\mathfrak{a},p\mathfrak{o})/p\mathfrak{o}$ であるからである。したがって、第一同型定理により、$\mathfrak{o}_p/\mathfrak{a}_p$ は、仮定により $\mathfrak{I}_p$ とともに、$\mathfrak{o}/(\mathfrak{a},p\mathfrak{o})$ と同型になる。

これは次を意味する。$\mathfrak{I}$ において $H(f_i(x))\equiv0\pmod p$ なら、$\mathfrak{o}$ において $H(u)\equiv0\pmod{(\mathfrak{a},p\mathfrak{o})}$ である。

これは次のように言い換えられる。$H(f(x))=p\cdot g(x)$（$H(f(x))$ は $\mathfrak{I}$ に、$g(x)$ は $\mathfrak{Z}$ に属する）なら、$H(u)-pF(u)\equiv0\pmod{\mathfrak{a}}$ が従い、したがって $H(f_i(x))=pF(f_i(x))$ であり、$g(x)=F(f_i(x))$ となる。ゆえに $g(x)$ は $\mathfrak{I}$ に属する。

有限回これを繰り返すと、$a$ がどの例外素数にも割られないなら、$a\cdot g(x)$ が $\mathfrak{I}$ に属し、$g(x)$ が $\mathfrak{Z}$ に属することから、常に $g(x)$ が $\mathfrak{I}$ に属することが従う。これで系は証明された。

\subsection*{\S\,3. 定理の証明}

\noindent\textbf{1. 補題.}
\emph{$\mathfrak{I}$ と $\mathfrak{o}$ の間の準同型を媒介する素イデアル $\mathfrak{a}$（\S\,2.3）は絶対素イデアルである。}

$\mathfrak{o}^*$ を、$u_1,\ldots,u_s$ に関する多項式領域であって、任意の代数数 --分数を含む-- を係数として許すものとする。また $\mathfrak{a}^*$ を、$\mathfrak{a}$ によって $\mathfrak{o}^*$ 内に生成されるイデアルとする。すなわち $\mathfrak{a}^*$ は、$\alpha_i$ が任意の代数数、$G_i(u)$ が $\mathfrak{a}$ の多項式であるような線型結合 $\alpha_1G_1(u)+\cdots+\alpha_sG_s(u)$ 全体から成る。証明すべきことは、$\mathfrak{a}^*$ が素イデアルであることである。

これは、$\mathfrak{a}^*$ が $\mathfrak{o}^*$ から $\mathfrak{I}^*$ への準同型を媒介することを示せばよい。ここで $\mathfrak{I}^*$ とは、$f_i(x)$ の多項式で係数として任意の代数数を許すもの全体からなる整域をいう。すなわち、$H(u)$ が $\mathfrak{o}^*$ に属し、$H(f_i(x))$ が消えるなら、$H(u)$ は $\mathfrak{a}^*$ に属することを示せばよい。

$H(u)$ の有限個の係数が生成する有限数体について、$\alpha_1,\ldots,\alpha_s$ を線型独立な基底とする。すると
\[
\alpha\cdot H(u)=\alpha_1H_1(u)+\cdots+\alpha_sH_s(u)
\]
となる。ただし $H_i(u)$ は $\mathfrak{o}$ に属し、$\alpha\ne0$ は有理整数である。$H(f)=0$ から
$\alpha_1H_1(f)+\cdots+\alpha_sH_s(f)=0$ が従い、したがって $H_1(f)=0,\ldots,H_s(f)=0$ となる。よって $H_1(u),\ldots,H_s(u)$ は $\mathfrak{a}$ に属し、したがって $H(u)$ は $\mathfrak{a}^*$ に属する。

\medskip\noindent\textbf{2.}
定理の証明は、いまや次の形で行える。$p$ を次のように選んだとする。
\begin{enumerate}
\item[I.] $\mathfrak{I}_p$ が定義されている。
\item[II.] $\mathfrak{a}_p$ が素イデアル、しかも絶対素イデアルのままである。
\item[III.] $\mathfrak{I}_p$ の超越次数が $\mathfrak{I}$ のそれに比べて下がっていない。
\end{enumerate}
このとき $\mathfrak{o}_p$ から $\mathfrak{I}_p$ への準同型は $\mathfrak{a}_p$ によって媒介される。これらの条件は、有限個を超えない例外素数を除き、すべての素数について満たされる。

条件 I, II, III に関して例外素数が有限個だけであることは、前述から直ちに従う。I は明らかである。例外素数は $f_i(x)$ の分母に入る有限個だけである。II は、補題により $\mathfrak{a}$ が絶対素イデアルであり、この性質を有限個を超えない素数を除くすべての素数を法として保つことから従う（注6参照）。最後に III は \S\,1.2 の系で示されている。なぜなら、$f_1(x),\ldots,f_t(x)$ の関数行列は \S\,2.2 により正確に階数 $t$ だからである。

いま $p$ をこれらの例外素数とは異なるように選ぶ。$\mathfrak{o}_p$ から $\mathfrak{I}_p$ への準同型は、あるイデアル $\mathfrak{c}_p$ によって媒介される。このイデアルは素イデアルである。なぜなら $\mathfrak{I}_p$ も零因子をもたない環だからである。またその超越次数は $\mathfrak{I}_p$ の超越次数によって与えられ、したがって少なくとも $t$ である。実際、$\mathfrak{I}_p$ は $\mathfrak{o}_p/\mathfrak{c}_p$ と同型であるから、$\mathfrak{c}_p$ の剰余類体は $\mathfrak{I}_p$ の商体と同型になる。

さらに $\mathfrak{a}_p$ も素イデアルで、$\mathfrak{c}_p$ の倍イデアル（部分イデアル）であるから、$\mathfrak{a}_p$ と $\mathfrak{c}_p$ は、その超越次数が一致するとき、かつそのときに限り一致する。\footnote{例えば B. L. v. d. Waerden, ``Zur Nullstellentheorie der Polynomideale'', Math. Ann. 96 (1926), pp. 183--208, \S\,3 を参照。} $\mathfrak{c}_p$ の倍イデアルとしての $\mathfrak{a}_p$ の超越次数も少なくとも $t$ であり、しかも $u_1,\ldots,u_t$ の $\mathfrak{a}_p$ を法とする類は既約系をなす。この超越次数がちょうど $t$ であることを示せば、すべてが証明される。

ここで、冪級数またはその商の場合には、$\mathfrak{I}$ の商体 $\mathfrak{Q}$ の構造に関する付加仮定（\S\,2.2）を用いる。それによれば、素イデアル $\mathfrak{a}$ はいずれの場合にも超越次数 $t$ であった。しかも番号付けは、$u_1,\ldots,u_t$ の $\mathfrak{a}$ を法とする類が既約系をなし、$u_{t+1},\ldots,u_s$ の類がそれらに代数的に依存するように選ばれていた。したがって $\mathfrak{a}$ には、整係数かつ原始的な多項式
\[
G_1(u),\ldots,G_{s-t}(u)\tag{1}
\]
が含まれており、これらは、原始性により消えない多項式
\[
\bar G_1(u),\ldots,\bar G_{s-t}(u)\tag{2}
\]
として $\mathfrak{a}_p$ に移る。しかし各 $\bar G_j$ には $u_{t+j}$ が実際に現れなければならない。そうでなければ、$p$ の選び方に反して、$u_1,\ldots,u_t$ の $\mathfrak{a}_p$ を法とする類の間に代数的依存関係が存在したことになるからである。こうして $\mathfrak{a}_p$ が超越次数 $t$ をもつことが示された。したがって $\mathfrak{a}_p$ は $\mathfrak{c}_p$ と同一であり、その準同型を媒介する。これで証明は完了する。

\subsection*{\S\,4. 全代数的係数への拡張}

導入部の終わりですでに述べたように、すべての考察は、有理整数の代わりに（有限な）代数的数体 $K$ の整数を、また素数 $p$ の代わりにその数体の素イデアル $\mathfrak{p}$ を基礎に置いても、わずかな修正で保たれる。このとき \S\,1 は完全にそのまま残り、\S\,2 における定理の定式化へ導く考察も同様である。ただし証明（\S\,3.2）および \S\,2.3 の系にはいくつかの補足が必要になる。

\medskip\noindent\textbf{1.}
\S\,3.2 では、さらに有限個の素イデアル、すなわち多項式 $G_1(u),\ldots,G_{s-t}(u)$（\S\,3.2, (1)）を割るものを例外素イデアルとして加えなければならない。というのも、いまやこれらを原始的と仮定することはできないからである。これによって $\bar G_1(u),\ldots,\bar G_{s-t}(u)$（\S\,3.2, (2)）が消えないことが保証される。これは条件 III の強化である。条件 II は保たれる。絶対既約性に関する定理は変わらないからである。もちろん I も保たれる。

\medskip\noindent\textbf{2.}
最終結果を表す \S\,2.3 の系は、そこでの最終形に対応して次のように定式化される。

\noindent\textbf{系.}
\emph{$K$ の整数 $a$ が有限個の例外素イデアルのいずれによっても割られないなら、$a\cdot g(x)$ が $\mathfrak{I}$ に属し、$g(x)$ が $\mathfrak{Z}$ に属することから、常に $g(x)$ が $\mathfrak{I}$ に属することが従う。}

$K$ の整数 $\lambda_1,\lambda_2,\ldots$ を、それぞれ一つの例外素イデアルで割られるが、その平方でも他の例外素イデアルでも割られないように選べば、分母 $a$ としてはこれらの $\lambda$ の冪積で十分である。

\noindent\emph{証明.}
$(a)=\mathfrak{p}_1^{e_1}\cdots\mathfrak{p}_k^{e_k}$ を $a$ の素イデアル分解とし、$\mathfrak{p}_1,\ldots,\mathfrak{p}_k$ は例外素イデアルではないとする。したがって各 $\mathfrak{p}_i$ について $\mathfrak{I}_{\mathfrak{p}_i}$ は $\mathfrak{o}_{\mathfrak{p}_i}/\mathfrak{a}_{\mathfrak{p}_i}$ と同型である。$K$ のすべての整数の環から、$(a)$ によって定められる商環 $R_a$ へ移る。これは、$(a)$ と素、すなわち $\mathfrak{p}_1,\ldots,\mathfrak{p}_k$ のいずれによっても割られない整数を、かつそれらだけを分母に加えることによって得られる。$R_a$ では $\mathfrak{p}_1,\ldots,\mathfrak{p}_k$ は素イデアルのままであるが主イデアルになり、他のすべての素イデアルは単位イデアルになる。また $(a)$ の分解は同じまま残る。\footnote{例えば H. Grell, ``Zur Theorie der Ordnungen in algebraischen Zahl- und Funktionenkörpern'', Math. Ann. 97 (1927), pp. 524--558, \S\,2.1 を参照。} したがって、係数領域を $R_a$ に移しても、系とその証明は \S\,2.3 の形そのままで保たれる。

つまり、$a\cdot g(x)$ が $\mathfrak{I}$ に属するなら、$g(x)=F(f_i(x))$ となり、ここで多項式 $F(u)$ の係数は $R_a$ に属する。これらの係数の共通分母を $d$ とすれば、$d$ は $a$ と素である。したがって次の形にも言い換えられる。$a\cdot g(x)$ が $\mathfrak{I}$ に属すれば、$a$ と素な $d$ について $d\cdot g(x)$ が $\mathfrak{I}$ に属する。ゆえに $(d,a)$ は単位イデアルであり、その結果、$a\cdot g(x)$ が $\mathfrak{I}$ に属することから $g(x)$ が $\mathfrak{I}$ に属することが従う。これで上の系の第一主張が証明された。

第二主張を証明するため、$R^*$ を有限個の例外素イデアルによって $K$ 内に定められる商環とし、$\lambda_1,\lambda_2,\ldots$ を $R^*$ におけるこれらの素イデアルの基底元とする。これらは整数として取ってよい。$R^*$ では、$a$ は $R^*$ の単元を除いて $\lambda$ の冪積になる。整数へ戻ると
\[
\gamma a=\beta\lambda_1^{\alpha_1}\cdots\lambda_r^{\alpha_r}
\]
が得られる。ただし $\gamma$ と $\beta$ は、どの例外素イデアルによっても割られない。$a\cdot g(x)$ が $\mathfrak{I}$ に属するなら、$\beta\lambda_1^{\alpha_1}\cdots\lambda_r^{\alpha_r}g(x)$ が $\mathfrak{I}$ に属し、さらに系の第一主張により $\lambda_1^{\alpha_1}\cdots\lambda_r^{\alpha_r}g(x)$ も $\mathfrak{I}$ に属する。

\begin{center}
（1931年11月11日受理）
\end{center}

\bigskip
\noindent\rule{\textwidth}{0.4pt}
\medskip
\noindent\textbf{英語要旨}（原文のロシア語要旨からの翻訳）:

\smallskip
\emph{任意の整域 $\mathfrak{I}$（Integritätsbereich、零因子をもたない環）は、同じ種類のある極大領域 $M(\mathfrak{I})$ を定める。それは、ある $a\ne0$ について $a\cdot g(x)$ が $\mathfrak{I}$ に含まれるようなすべての関数 $g(x)$ から成る。}

\emph{本論文では、この極大領域 $M(\mathfrak{I})$ の「分母」$a$ の問題を、領域 $\mathfrak{I}$ が有限であったという仮定のもとで考察する。本論文で証明されるのは、これらの $a$ を、そこに因子として現れる素数が全体として有限個だけになるように常に選べる、ということである。}

\emph{同時に、$g(x)$ が冪級数またはそのような級数の商である場合には、領域 $\mathfrak{I}$ を定義する関数 $f_i(x)$ が相互に代数関係だけで結ばれ得るものと仮定する。}

\emph{証明は、問題を一般イデアル論のいくつかの問題へ帰着することに基づく。この帰着は、分母 $a$ に素数 $p$ が現れることを、関数 $f_i$ の間の法 $p$ の関係と見なすことによって行われる。}
